% -----------------------------------------------------------------------------
% Chapter 1 LaTeX version
% Source QMD: chapter-01-from-single-variable-calculus-to-multivariable-thinking.qmd
% This file is self-contained. To use inside a larger book, copy the material
% after \begin{document}, or move the preamble definitions to the main preamble.
% -----------------------------------------------------------------------------

\documentclass[11pt,oneside]{book}

\usepackage[margin=1in]{geometry}
\usepackage[T1]{fontenc}
\usepackage[utf8]{inputenc}
\usepackage{lmodern}
\usepackage{microtype}
\usepackage{amsmath,amssymb,mathtools,bm}
\usepackage{booktabs,array,tabularx,multirow}
\usepackage{enumitem}
\usepackage{hyperref}
\usepackage{xcolor}
\usepackage{tikz}
\usepackage{pgfplots}
\usepackage{float}
\usepackage{caption}
\usepackage{listings}
\usepackage[most]{tcolorbox}

\pgfplotsset{compat=1.17}
\usetikzlibrary{arrows.meta,calc,positioning,shadows.blur,decorations.pathreplacing,patterns,backgrounds,fit,intersections}

% -----------------------------------------------------------------------------
% Colors and style
% -----------------------------------------------------------------------------
\definecolor{chapterblue}{HTML}{1F5AA6}
\definecolor{chapterteal}{HTML}{168A8F}
\definecolor{chapterorange}{HTML}{D9822B}
\definecolor{chaptergreen}{HTML}{2D8A4E}
\definecolor{chapterpurple}{HTML}{6B4EAA}
\definecolor{softblue}{HTML}{EAF3FF}
\definecolor{softgreen}{HTML}{EAF8EF}
\definecolor{softorange}{HTML}{FFF4E6}
\definecolor{softpurple}{HTML}{F3EEFF}
\definecolor{softgray}{HTML}{F6F7F9}
\definecolor{darkgray}{HTML}{30343B}

\hypersetup{
  colorlinks=true,
  linkcolor=chapterblue,
  urlcolor=chapterteal,
  citecolor=chapterpurple
}

\setlength{\parskip}{0.45em}
\setlength{\parindent}{0pt}
\setlist[itemize]{topsep=0.3em,itemsep=0.2em}
\setlist[enumerate]{topsep=0.3em,itemsep=0.2em}

\captionsetup{font=small,labelfont=bf}

\tikzset{
  axisline/.style={-Latex,thick,darkgray},
  vector/.style={-Latex,very thick,chapterblue},
  smallvector/.style={-Latex,thick,chapterblue},
  dashedline/.style={dash pattern=on 4pt off 2pt,gray},
  point/.style={circle,fill=chapterorange,inner sep=1.5pt},
  boxnode/.style={rounded corners=3pt,draw=chapterblue,fill=softblue,thick,align=center,inner sep=6pt},
  greennode/.style={rounded corners=3pt,draw=chaptergreen,fill=softgreen,thick,align=center,inner sep=6pt},
  orangenode/.style={rounded corners=3pt,draw=chapterorange,fill=softorange,thick,align=center,inner sep=6pt},
  purplenode/.style={rounded corners=3pt,draw=chapterpurple,fill=softpurple,thick,align=center,inner sep=6pt}
}

\newtcolorbox{mainideabox}[1][] {
  enhanced, breakable,
  colback=softblue,
  colframe=chapterblue,
  coltitle=white,
  fonttitle=\bfseries,
  title={Main idea},
  sharp corners=south,
  boxrule=0.8pt,
  left=1mm,right=1mm,top=1mm,bottom=1mm,
  #1
}

\newtcolorbox{bridgebox}[1][] {
  enhanced, breakable,
  colback=softgreen,
  colframe=chaptergreen,
  fonttitle=\bfseries,
  title={Bridge},
  boxrule=0.8pt,
  left=1mm,right=1mm,top=1mm,bottom=1mm,
  #1
}

\newtcolorbox{definitionbox}[1][] {
  enhanced, breakable,
  colback=softpurple,
  colframe=chapterpurple,
  fonttitle=\bfseries,
  title={Definition},
  boxrule=0.8pt,
  left=1mm,right=1mm,top=1mm,bottom=1mm,
  #1
}

\newtcolorbox{examplebox}[1][] {
  enhanced, breakable,
  colback=softorange,
  colframe=chapterorange,
  fonttitle=\bfseries,
  title={Example},
  boxrule=0.8pt,
  left=1mm,right=1mm,top=1mm,bottom=1mm,
  #1
}

\newtcolorbox{warningbox}[1][] {
  enhanced, breakable,
  colback=softgray,
  colframe=darkgray,
  fonttitle=\bfseries,
  title={Common misconception},
  boxrule=0.8pt,
  left=1mm,right=1mm,top=1mm,bottom=1mm,
  #1
}

\lstdefinestyle{pythonstyle}{
  language=Python,
  basicstyle=\ttfamily\small,
  keywordstyle=\color{chapterblue}\bfseries,
  stringstyle=\color{chaptergreen},
  commentstyle=\color{gray},
  numberstyle=\tiny\color{gray},
  numbers=left,
  stepnumber=1,
  numbersep=7pt,
  showstringspaces=false,
  breaklines=true,
  frame=single,
  rulecolor=\color{gray!50},
  backgroundcolor=\color{softgray}
}

\newcommand{\R}{\mathbb{R}}
\newcommand{\grad}{\nabla}
\newcommand{\dd}{\,d}
\newcommand{\vect}[1]{\left\langle #1\right\rangle}

\begin{document}
\setcounter{chapter}{0}

\chapter{From Single-Variable Calculus to Multivariable Thinking}
\label{ch:single-to-multivariable}

\section*{Chapter overview}

Single-variable calculus studies functions such as $f(x)$, where one input number produces one output number. Multivariable calculus studies functions whose inputs may be points, vectors, measurements, parameters, or data records:
\[
  f:\R^n \longrightarrow \R,
  \qquad
  x=(x_1,x_2,\dots,x_n) \longmapsto f(x).
\]

This change looks small, but it changes the whole subject. In one-variable calculus, we move left or right along a line. In multivariable calculus, we can move in infinitely many directions. In one-variable calculus, the graph of $y=f(x)$ is a curve. In two-variable calculus, the graph of $z=f(x,y)$ is a surface. In higher dimensions, the graph may be impossible to draw directly, but the same ideas remain: local approximation, rate of change, accumulation, optimization, and geometric structure.

This chapter is a bridge. It reviews core ideas of single-variable calculus and explains how they become the language of multivariable calculus, applied mathematics, statistics, data science, and machine learning.

\begin{mainideabox}
Multivariable calculus is the calculus of \textbf{many inputs}. The output may represent temperature, height, cost, probability density, potential energy, risk, error, or loss. The input may represent spatial coordinates, model parameters, data features, or the state of a system.
\end{mainideabox}

\begin{figure}[H]
\centering
\begin{tikzpicture}[scale=0.95]
  % Left: single-variable
  \node[boxnode,minimum width=3.1cm,minimum height=1.1cm] (single) at (0,0) {$f:\R\to\R$\\ one input};
  \draw[axisline] (-3.2,-1.7) -- (3.2,-1.7) node[right] {$x$};
  \draw[chapterblue,very thick,smooth,domain=-2.4:2.4,samples=40] plot (\x,{0.18*(\x)^3-0.7*(\x)-1.7});
  \node[darkgray] at (0,-3.0) {curve $y=f(x)$};

  % Arrow
  \draw[-Latex,very thick,chapterorange] (3.8,-0.7) -- (5.2,-0.7) node[midway,above] {many inputs};

  % Right: multivariable
  \node[greennode,minimum width=3.5cm,minimum height=1.1cm] (multi) at (8.3,0) {$f:\R^2\to\R$\\ two inputs};
  \begin{scope}[shift={(8.3,-2.0)},x={(0.85cm,-0.23cm)},y={(0.95cm,0.12cm)},z={(0cm,0.75cm)}]
    \draw[axisline] (0,0,0) -- (2.4,0,0) node[below] {$x$};
    \draw[axisline] (0,0,0) -- (0,2.2,0) node[right] {$y$};
    \draw[axisline] (0,0,0) -- (0,0,2.4) node[left] {$z$};
    \foreach \r in {0.35,0.7,1.05,1.4} {
      \draw[chapterteal!70] plot[smooth,domain=0:360,samples=80] ({\r*cos(\x)},{\r*sin(\x)},0);
    }
    \draw[chaptergreen,very thick] plot[smooth,domain=-1.5:1.5,samples=80] (\x,{0},{0.25*(\x)^2+0.45});
    \node[chaptergreen] at (1.4,0,1.9) {surface};
  \end{scope}
\end{tikzpicture}
\caption{The central shift: from a one-dimensional input line to a multidimensional input space.}
\label{fig:one-input-to-many-inputs}
\end{figure}

\section*{Learning goals}

After studying this chapter, you should be able to:
\begin{enumerate}
  \item Explain the difference between $f:\R\to\R$, $f:\R^2\to\R$, $f:\R^3\to\R$, and $f:\R^n\to\R$.
  \item Interpret a multivariable input as a point, vector, data record, parameter vector, or state.
  \item Describe a function using formulas, tables, graphs, level sets, and computation.
  \item Explain why derivatives become directional and why gradients are central.
  \item Explain why integrals become accumulation over regions, volumes, and probability spaces.
  \item Recognize the role of multivariable calculus in optimization, statistics, machine learning, and applied mathematics.
\end{enumerate}

\section{Why multivariable calculus?}

In single-variable calculus, a typical function has the form
\[
  f:\R\to\R,
  \qquad
  x\mapsto f(x).
\]
For example, $f(x)=x^2$ takes one number $x$ and returns one number $x^2$. This is already powerful. We can ask for slopes, tangent lines, areas under curves, maximum values, and differential equations.

But modern quantitative problems rarely depend on only one variable. A temperature model may depend on longitude, latitude, altitude, and time:
\[
  T=T(x,y,z,t).
\]
A statistical loss function may depend on many model parameters:
\[
  L=L(\beta_0,\beta_1,\dots,\beta_p).
\]
A neural network may have millions or billions of parameters, and training means minimizing a function of all those parameters:
\[
  \text{loss}=L(w_1,w_2,\dots,w_N).
\]
A probability density for two continuous random variables is a function
\[
  p:\R^2\to\R,
  \qquad
  (x,y)\mapsto p(x,y),
\]
and a velocity field assigns a vector to each point in space:
\[
  F:\R^3\to\R^3.
\]

So the core move is
\[
  \text{one input} \qquad \longrightarrow \qquad \text{many inputs}.
\]

\begin{bridgebox}[title={Bridge to MS preparation}]
For applied mathematics, statistics, data science, AI, and machine learning, the most important mental shift is to see formulas as functions on spaces. A loss function is a surface or landscape. A density is a surface or high-dimensional object whose total mass is one. A dynamical system is motion through a state space. A machine learning model is a map from features to predictions.
\end{bridgebox}

\section{Inputs, outputs, and spaces}

A function is a rule that assigns to each input exactly one output. The input set is the \textbf{domain}. The output set is the \textbf{codomain}. In calculus, the most common spaces are Euclidean spaces:
\[
  \R,\qquad \R^2,\qquad \R^3,\qquad \R^n.
\]
The space $\R^n$ consists of ordered $n$-tuples of real numbers:
\[
  \R^n=\{(x_1,x_2,\dots,x_n):x_i\in\R\}.
\]
The same $n$-tuple may be interpreted in different ways depending on the model.

\begin{center}
\begin{tabularx}{0.95\textwidth}{>{\raggedright\arraybackslash}p{0.28\textwidth} X}
\toprule
\textbf{Mathematical object} & \textbf{Possible interpretation}\\
\midrule
$(x,y)$ & A point in the plane, longitude and latitude, two features of an observation\\
$(x,y,z)$ & A point in space, three physical coordinates, RGB color values\\
$(x_1,\dots,x_n)$ & A data vector, a parameter vector, a state vector, a list of measurements\\
$f(x,y)$ & Height, temperature, cost, density, loss, utility\\
$F(x,y,z)$ & Vector field, force field, velocity field, gradient flow\\
\bottomrule
\end{tabularx}
\end{center}

\begin{definitionbox}[title={Scalar-valued and vector-valued functions}]
A \textbf{scalar-valued function} has real-number output:
\[
  f:\R^n\to\R.
\]
Examples include
\[
  f(x,y)=x^2+y^2,
  \qquad
  T(x,y,z)=x^2e^y+\ln z,
  \qquad
  L(w,b)=\frac1m\sum_{i=1}^m(wx_i+b-y_i)^2.
\]
A \textbf{vector-valued function} has vector output:
\[
  F:\R^n\to\R^m.
\]
Examples include
\[
  r(t)=\vect{\cos t,\sin t,t},
  \qquad
  F(x,y)=\vect{-y,x},
  \qquad
  G(x,y,z)=\vect{yz,xz,xy}.
\]
\end{definitionbox}

\begin{figure}[H]
\centering
\begin{tikzpicture}[node distance=1.1cm and 1.2cm]
  \node[boxnode,minimum width=3.2cm] (rr) {$f:\R\to\R$\\ curve};
  \node[greennode,minimum width=3.2cm,right=of rr] (r2r) {$f:\R^2\to\R$\\ surface};
  \node[purplenode,minimum width=3.2cm,right=of r2r] (r3r) {$f:\R^3\to\R$\\ level surfaces};
  \node[orangenode,minimum width=3.2cm,below=of r2r] (field) {$F:\R^n\to\R^m$\\ vector-valued map};
  \draw[-Latex,thick] (rr) -- (r2r);
  \draw[-Latex,thick] (r2r) -- (r3r);
  \draw[-Latex,thick] (r2r) -- (field);
  \node[align=center,below=0.9cm of field,text width=0.82\textwidth] {
  The output type determines the geometry: curves, surfaces, level sets, vector fields, transformations, and systems.};
\end{tikzpicture}
\caption{A small taxonomy of functions that appear in multivariable calculus.}
\label{fig:function-taxonomy}
\end{figure}

\section{A comparison with single-variable calculus}

The table below gives a roadmap for the course.

\begin{center}
\begin{tabularx}{\textwidth}{>{\raggedright\arraybackslash}p{0.25\textwidth} >{\raggedright\arraybackslash}p{0.31\textwidth} X}
\toprule
\textbf{Single-variable idea} & \textbf{Multivariable version} & \textbf{Modern interpretation}\\
\midrule
Number $x$ & Vector $x=(x_1,\dots,x_n)$ & Feature vector or parameter vector\\
Function $f(x)$ & Function $f(x_1,\dots,x_n)$ & Loss, density, energy, response surface\\
Graph $y=f(x)$ & Surface $z=f(x,y)$ or high-dimensional graph & Landscape\\
Slope $f'(a)$ & Partial derivatives and gradient $\grad f(a)$ & Sensitivity and learning signal\\
Tangent line & Tangent plane or linear approximation & Local linear model\\
Critical point $f'(a)=0$ & Critical point $\grad f(a)=0$ & Candidate optimum or saddle point\\
Integral $\int_a^b f(x)\dd x$ & Double or triple integral & Accumulation, mass, probability\\
Substitution & Change of variables and Jacobian & Transformation of densities and coordinates\\
Fundamental theorem of calculus & Green, Stokes, divergence theorem & Conservation and flow laws\\
\bottomrule
\end{tabularx}
\end{center}

The unifying theme is that calculus studies \textbf{change} and \textbf{accumulation}. Multivariable calculus studies change and accumulation when inputs live in spaces with dimension larger than one.

\section{Four ways to represent a function}

A modern calculus course should move comfortably among several representations:
\begin{enumerate}
  \item \textbf{Verbal}: a description in words.
  \item \textbf{Numerical}: a table, dataset, or sampled values.
  \item \textbf{Visual}: a graph, contour plot, heatmap, vector field, or surface.
  \item \textbf{Algebraic}: a formula.
  \item \textbf{Computational}: a function implemented in code.
\end{enumerate}

\subsection*{Example: temperature on a plate}

Suppose the temperature on a rectangular metal plate is modeled by
\[
  T(x,y)=80-2x^2-y^2.
\]
The same object can be represented in several ways.

\begin{center}
\begin{tabularx}{0.9\textwidth}{p{0.22\textwidth}X}
\toprule
\textbf{Representation} & \textbf{Description}\\
\midrule
Verbal & Temperature is highest near the origin and decreases quadratically away from it.\\
Algebraic & $T(x,y)=80-2x^2-y^2$\\
Visual & Surface plot or contour plot\\
Numerical & Table of values on a grid of $(x,y)$ points\\
Computational & A Python function \texttt{T(x, y)}\\
\bottomrule
\end{tabularx}
\end{center}

\begin{figure}[H]
\centering
\begin{tikzpicture}[scale=0.95]
  % coordinate axes
  \draw[axisline] (-3.8,0) -- (3.8,0) node[right] {$x$};
  \draw[axisline] (0,-3.1) -- (0,3.1) node[above] {$y$};
  % ellipses T=c, 2x^2+y^2 = 80-c
  \foreach \a/\b/\lab in {0.90/1.27/76,1.35/1.91/72,1.80/2.55/64,2.25/3.18/48} {
    \draw[chapterblue,very thick] (0,0) ellipse[x radius=\a,y radius=\b];
    \node[chapterblue,fill=white,inner sep=1pt] at ({\a*0.72},{\b*0.72}) {$T=\lab$};
  }
  \fill[chapterorange] (0,0) circle (2pt) node[above right] {maximum $80$};
  \node[align=center,text width=0.75\textwidth] at (0,-4.0) {Level curves of $T(x,y)=80-2x^2-y^2$. Each ellipse contains points with equal temperature.};
\end{tikzpicture}
\caption{Contour plot for the temperature model. The ellipses get larger as the temperature decreases.}
\label{fig:temperature-contours-tikz}
\end{figure}

The contour plot does not draw the full surface in three dimensions. Instead, it draws curves where the function has constant value. These curves are called \textbf{level curves}. They are among the most important visualization tools in multivariable calculus.

\section{Graphs, level curves, and level sets}

For a function of one variable $f:\R\to\R$, the graph is a subset of $\R^2$:
\[
  \{(x,y)\in\R^2:y=f(x)\}.
\]
For a function of two variables $f:\R^2\to\R$, the graph is a subset of $\R^3$:
\[
  \{(x,y,z)\in\R^3:z=f(x,y)\}.
\]
For example,
\[
  f(x,y)=x^2+y^2
\]
has graph
\[
  z=x^2+y^2,
\]
an elliptic paraboloid.

A \textbf{level curve} is a curve where the output is constant:
\[
  f(x,y)=c.
\]
For $f(x,y)=x^2+y^2$, the level curves are
\[
  x^2+y^2=c,
\]
which are circles centered at the origin when $c>0$.

\begin{figure}[H]
\centering
\begin{tikzpicture}
\begin{axis}[
  width=0.78\textwidth,
  height=0.48\textwidth,
  view={120}{30},
  xlabel={$x$},ylabel={$y$},zlabel={$z$},
  domain=-2.3:2.3,
  y domain=-2.3:2.3,
  samples=17,
  colormap/viridis,
  zmin=0,zmax=11,
  axis lines=center,
  ticklabel style={font=\scriptsize},
  label style={font=\small}
]
\addplot3[surf,shader=interp,opacity=0.82] {x^2+y^2};
\addplot3[domain=0:360,samples=40,very thick,chapterorange] ({1.0*cos(x)},{1.0*sin(x)},{0});
\addplot3[domain=0:360,samples=40,very thick,chapterorange] ({1.6*cos(x)},{1.6*sin(x)},{0});
\addplot3[domain=0:360,samples=40,very thick,chapterorange] ({2.1*cos(x)},{2.1*sin(x)},{0});
\end{axis}
\end{tikzpicture}
\caption{The surface $z=x^2+y^2$ together with circular level curves projected onto the $xy$-plane.}
\label{fig:paraboloid-with-level-curves}
\end{figure}

\begin{mainideabox}[title={Why level sets matter}]
In higher dimensions, full graphs are usually impossible to draw. Level sets remain useful. For $f:\R^n\to\R$, a level set has the form
\[
  \{x\in\R^n:f(x)=c\}.
\]
In statistics, level sets of a likelihood or log-likelihood show parameter values with equal fit. In machine learning, level sets of a loss function show parameter values with equal training error.
\end{mainideabox}

\section{Derivatives become directional}

In one-variable calculus, the derivative is the limit
\[
  f'(a)=\lim_{h\to0}\frac{f(a+h)-f(a)}{h}.
\]
This measures the instantaneous rate of change as $x$ moves along the real line.

For a function $f(x,y)$, the input point $(x,y)$ can move in many directions. From $(a,b)$, it can move in the $x$ direction, the $y$ direction, diagonally, along a curve, or along any unit vector $u$.

This leads to several related ideas:
\begin{itemize}
  \item \textbf{Partial derivatives}: rates of change along coordinate directions.
  \item \textbf{Directional derivatives}: rates of change along arbitrary directions.
  \item \textbf{Gradient}: the vector that organizes all directional derivative information.
  \item \textbf{Jacobian}: the matrix that represents the derivative of a vector-valued function.
\end{itemize}

For a scalar-valued function $f(x,y)$, the gradient is
\[
  \grad f(x,y)=\left\langle \frac{\partial f}{\partial x}(x,y),\frac{\partial f}{\partial y}(x,y)\right\rangle.
\]
For $f(x,y,z)$,
\[
  \grad f(x,y,z)=\left\langle \frac{\partial f}{\partial x},\frac{\partial f}{\partial y},\frac{\partial f}{\partial z}\right\rangle.
\]
The gradient points in the direction of fastest increase. This is why gradients are central in optimization and machine learning.

\begin{examplebox}[title={A first gradient computation}]
Let
\[
  f(x,y)=x^2+2y^2.
\]
Then
\[
  \frac{\partial f}{\partial x}=2x,
  \qquad
  \frac{\partial f}{\partial y}=4y,
\]
so
\[
  \grad f(x,y)=\vect{2x,4y}.
\]
At the point $(1,1)$,
\[
  \grad f(1,1)=\vect{2,4}.
\]
Thus near $(1,1)$, the function increases fastest in the direction of $\vect{2,4}$.
\end{examplebox}

\begin{figure}[H]
\centering
\begin{tikzpicture}[scale=1.05]
  \draw[axisline] (-3.2,0) -- (3.2,0) node[right] {$x$};
  \draw[axisline] (0,-2.6) -- (0,2.6) node[above] {$y$};
  % level curves x^2+2y^2=c
  \foreach \a/\b/\lab in {0.8/0.57/1,1.35/0.95/3,1.9/1.34/6,2.45/1.73/10} {
    \draw[chapterteal,very thick] (0,0) ellipse[x radius=\a,y radius=\b];
  }
  % gradient arrows along radial-ish directions
  \foreach \x/\y in {-2/-1,-1.4/1.1,-0.7/-1.4,0.8/1.2,1.5/-1.1,2/0.8} {
    \pgfmathsetmacro{\u}{0.35*2*\x}
    \pgfmathsetmacro{\v}{0.35*4*\y}
    \draw[smallvector,chapterorange] (\x,\y) -- ++(\u,\v);
  }
  \fill[chapterorange] (1,1) circle (2pt) node[above right] {$(1,1)$};
  \draw[vector,chapterpurple] (1,1) -- ++(0.65,1.3) node[right] {$\grad f(1,1)=\vect{2,4}$};
  \node[align=center,text width=0.82\textwidth] at (0,-3.35) {For $f(x,y)=x^2+2y^2$, gradients are perpendicular to the elliptical level curves and point toward larger function values.};
\end{tikzpicture}
\caption{Gradient vectors organize directional change.}
\label{fig:gradient-on-contours-tikz}
\end{figure}

\section{Integrals become accumulation over regions}

In one-variable calculus, the definite integral
\[
  \int_a^b f(x)\dd x
\]
accumulates values of $f$ over an interval. In multivariable calculus, we accumulate over regions and solids:
\[
  \iint_R f(x,y)\dd A,
  \qquad
  \iiint_E f(x,y,z)\dd V.
\]
If $f(x,y)\ge0$, then
\[
  \iint_R f(x,y)\dd A
\]
can be interpreted as the volume under the surface $z=f(x,y)$ above the region $R$.

If $\rho(x,y,z)$ is a density, then
\[
  \iiint_E \rho(x,y,z)\dd V
\]
is the mass of the object occupying the solid region $E$.

If $p(x,y)$ is a joint probability density, then
\[
  \iint_R p(x,y)\dd A
\]
is the probability that $(X,Y)$ lies in the region $R$.

\begin{figure}[H]
\centering
\begin{tikzpicture}[scale=0.92, x={(1cm,0cm)}, y={(0.55cm,0.28cm)}, z={(0cm,0.88cm)}]
  % base grid
  \foreach \i in {0,...,5} {
    \draw[gray!35] (\i,0,0) -- (\i,4,0);
  }
  \foreach \j in {0,...,4} {
    \draw[gray!35] (0,\j,0) -- (5,\j,0);
  }
  \draw[axisline] (0,0,0) -- (5.6,0,0) node[below] {$x$};
  \draw[axisline] (0,0,0) -- (0,4.7,0) node[right] {$y$};
  \draw[axisline] (0,0,0) -- (0,0,4.1) node[left] {$z$};
  \fill[chapterorange!25,opacity=0.8] (0,0,0) -- (5,0,0) -- (5,4,0) -- (0,4,0) -- cycle;
  \node at (2.5,2,0) {$R$};
  % columns
  \foreach \i/\j/\h in {0.7/0.6/1.1,1.7/0.6/1.7,2.7/0.6/2.1,3.7/0.6/2.4,0.7/1.6/1.5,1.7/1.6/2.0,2.7/1.6/2.6,3.7/1.6/2.8,0.7/2.6/1.7,1.7/2.6/2.3,2.7/2.6/3.0,3.7/2.6/3.2} {
    \draw[chapterblue!60,fill=chapterblue!16] (\i,\j,0) -- ++(0.5,0,0) -- ++(0,0.5,0) -- ++(-0.5,0,0) -- cycle;
    \draw[chapterblue!60] (\i,\j,0) -- (\i,\j,\h);
    \draw[chapterblue!60] (\i+0.5,\j,0) -- (\i+0.5,\j,\h);
    \draw[chapterblue!60] (\i+0.5,\j+0.5,0) -- (\i+0.5,\j+0.5,\h);
    \draw[chapterblue!60] (\i,\j+0.5,0) -- (\i,\j+0.5,\h);
    \fill[chapterblue!30,opacity=0.85] (\i,\j,\h) -- ++(0.5,0,0) -- ++(0,0.5,0) -- ++(-0.5,0,0) -- cycle;
  }
  % surface sketch
  \draw[chaptergreen,very thick,smooth] plot coordinates {(0.5,0.2,1.0) (1.5,0.7,1.8) (2.6,1.3,2.5) (3.8,2.0,3.1) (4.8,3.3,3.5)};
  \node[chaptergreen] at (5.2,3.6,3.5) {$z=f(x,y)$};
  \node[align=center,text width=0.76\textwidth] at (2.4,-1.1,0) {A double integral adds many small columns: $\iint_R f(x,y)\dd A$.};
\end{tikzpicture}
\caption{A double integral as accumulated volume above a region.}
\label{fig:double-integral-volume-tikz}
\end{figure}

\begin{mainideabox}[title={Statistics preview}]
For a joint density $p(x,y)$, total probability means
\[
  \iint_{\R^2}p(x,y)\dd A=1.
\]
Expected values, variances, covariances, and conditional distributions are all built from multivariable integrals.
\end{mainideabox}

\section{The geometry of optimization}

Optimization is one of the main bridges from calculus to statistics and machine learning. In one variable, we minimize or maximize a curve $f(x)$. In several variables, we minimize or maximize a surface or high-dimensional landscape $f(x_1,\dots,x_n)$.

A basic unconstrained optimization problem is
\[
  \min_{x\in\R^n} f(x).
\]
For a differentiable function, candidate optima often satisfy
\[
  \grad f(x)=0.
\]
In machine learning, the function $f$ is often a \textbf{loss function}. Training a model means finding parameters that make the loss small.

\subsection*{Example: least-squares loss for a line}

Suppose we have data points $(x_i,y_i)$ and want to fit a line
\[
  \widehat y=wx+b.
\]
The mean squared error is
\[
  L(w,b)=\frac1m\sum_{i=1}^m(wx_i+b-y_i)^2.
\]
This is a function of two variables, $w$ and $b$. The input point $(w,b)$ is not a physical point in space; it is a point in \textbf{parameter space}.

\begin{figure}[H]
\centering
\begin{tikzpicture}[scale=0.95]
  % left contour landscape
  \begin{scope}[shift={(0,0)}]
    \draw[axisline] (-3.2,0) -- (3.2,0) node[right] {$w$};
    \draw[axisline] (0,-2.6) -- (0,2.6) node[above] {$b$};
    \foreach \a/\b in {0.6/0.35,1.1/0.65,1.7/1.0,2.35/1.38,2.9/1.7} {
      \draw[chapterblue,very thick,rotate=20] (0.5,-0.3) ellipse[x radius=\a,y radius=\b];
    }
    \fill[chapterorange] (0.5,-0.3) circle (2pt) node[below right] {minimum};
    \draw[chapterorange,very thick,-Latex] (-2.4,1.8) -- (-1.4,1.1) -- (-0.65,0.55) -- (-0.1,0.15) -- (0.28,-0.15) -- (0.5,-0.3);
    \node[align=center] at (0,-3.15) {contours of $L(w,b)$};
  \end{scope}
  % right 3D bowl
  \begin{scope}[shift={(7.3,-0.1)}]
    \begin{axis}[
      width=0.46\textwidth,
      height=0.34\textwidth,
      view={120}{32},
      xlabel={$w$},ylabel={$b$},zlabel={$L$},
      domain=-2:2,
      y domain=-2:2,
      samples=17,
      colormap/viridis,
      axis lines=center,
      ticklabel style={font=\scriptsize},
      label style={font=\small}
    ]
      \addplot3[surf,shader=interp,opacity=0.82] {(x-0.4)^2+2*(y+0.2)^2+0.4*x*y};
    \end{axis}
  \end{scope}
\end{tikzpicture}
\caption{A least-squares loss surface can be viewed as a bowl-shaped landscape in parameter space.}
\label{fig:loss-landscape-tikz}
\end{figure}

The goal is to move downhill on this surface. In later chapters, gradient descent will update parameters by
\[
  (w,b)\leftarrow (w,b)-\eta\grad L(w,b),
\]
where $\eta>0$ is a step size called the learning rate.

\begin{bridgebox}[title={AI connection}]
Training a neural network is conceptually similar to minimizing $L(w,b)$, but the parameter vector may have millions of components. The derivative is not just a slope. It is a high-dimensional gradient, and the chain rule becomes backpropagation.
\end{bridgebox}

\section{Local and global viewpoints}

Calculus has both local and global ideas.

A \textbf{local} idea describes what happens near one point. Examples include tangent lines, tangent planes, derivatives, gradients, linear approximation, Hessian matrices, local maxima, local minima, and saddle points.

A \textbf{global} idea describes an entire interval, region, surface, or space. Examples include total area, total volume, total mass, total probability, line integrals around closed curves, flux through a surface, and global maximum and minimum values.

\begin{figure}[H]
\centering
\begin{tikzpicture}
  \begin{axis}[
    width=0.78\textwidth,
    height=0.46\textwidth,
    view={125}{32},
    xlabel={$x$},ylabel={$y$},zlabel={$z$},
    domain=-1.8:1.8,
    y domain=-1.8:1.8,
    samples=17,
    colormap/viridis,
    axis lines=center,
    ticklabel style={font=\scriptsize},
    label style={font=\small},
    zmin=-1,zmax=8
  ]
    \addplot3[surf,shader=interp,opacity=0.75] {x^2+0.7*y^2+0.2*x*y+1};
    \addplot3[surf,domain=-0.8:0.8,y domain=-0.8:0.8,samples=2,opacity=0.55,fill=chapterorange,draw=chapterorange] {2 + 2*x + 1.2*y};
    \addplot3[only marks,mark=*,mark size=2.5pt,chapterorange] coordinates {(0,0,2)};
  \end{axis}
\end{tikzpicture}
\caption{A local tangent plane approximates a surface near one point, while global questions concern the whole surface or region.}
\label{fig:local-tangent-plane-global-surface}
\end{figure}

Both viewpoints are essential. A gradient is local, but gradient descent uses many local steps to solve a global problem. A density may be defined pointwise, but probabilities require integrating the density over a region. A vector field may assign a vector at each point, but flux and circulation describe global behavior.

\section{High-dimensional thinking}

Most of this book builds geometric intuition in $\R^2$ and $\R^3$. But the real goal is to understand ideas that work in $\R^n$.

A data point with $n$ features can be written as
\[
  x=(x_1,x_2,\dots,x_n)\in\R^n.
\]
A linear prediction model may have the form
\[
  \widehat y=\beta_0+\beta_1x_1+\cdots+\beta_nx_n.
\]
A loss function for the parameters $\beta=(\beta_0,\beta_1,\dots,\beta_n)$ may be written as
\[
  L(\beta)=\frac1m\sum_{i=1}^m(\beta_0+\beta_1x_{i1}+\cdots+\beta_nx_{in}-y_i)^2.
\]
This is a scalar-valued function on a high-dimensional parameter space. Even when we cannot draw the graph, we can still compute
\[
  \grad L(\beta),
  \qquad
  \grad^2 L(\beta),
  \qquad
  \min_\beta L(\beta).
\]
These are multivariable calculus objects.

\begin{figure}[H]
\centering
\begin{tikzpicture}[node distance=0.6cm]
  \node[boxnode,minimum width=8.4cm] (x) {$x=(x_1,x_2,\dots,x_n)$\\ feature vector};
  \node[purplenode,below=of x,minimum width=8.4cm] (beta) {$\beta=(\beta_0,\beta_1,\dots,\beta_n)$\\ parameter vector};
  \node[orangenode,below=of beta,minimum width=8.4cm] (model) {$\widehat y=\beta_0+\beta_1x_1+\cdots+\beta_nx_n$\\ prediction};
  \node[greennode,below=of model,minimum width=8.4cm] (loss) {$L(\beta)=\dfrac1m\sum(\widehat y_i-y_i)^2$\\ loss landscape in parameter space};
  \draw[-Latex,thick] (x) -- (beta);
  \draw[-Latex,thick] (beta) -- (model);
  \draw[-Latex,thick] (model) -- (loss);
  \node[align=center,text width=0.75\textwidth,below=0.55cm of loss] {High-dimensional calculus uses the same objects: vectors, functions, gradients, Hessians, integrals, and approximations.};
\end{tikzpicture}
\caption{A high-dimensional model pipeline: features, parameters, predictions, and loss.}
\label{fig:high-dimensional-pipeline}
\end{figure}

\begin{warningbox}
High-dimensional calculus is not a different subject. It is the same subject with fewer pictures and more structure. We use lower-dimensional pictures to build intuition, then rely on vectors, matrices, gradients, Jacobians, Hessians, and integrals to work in general dimensions.
\end{warningbox}

\section{Worked examples}

\begin{examplebox}[title={1. Classifying functions by domain and output}]
Classify each function by its domain and output type.
\begin{enumerate}
  \item $f(x)=x^2+1$.
  \item $g(x,y)=\sqrt{x^2+y^2}$.
  \item $T(x,y,z)=x^2e^y+\ln z$.
  \item $r(t)=\vect{\cos t,\sin t,t}$.
  \item $F(x,y)=\vect{-y,x}$.
\end{enumerate}

\textbf{Solution.}
\begin{enumerate}
  \item $f:\R\to\R$ is a scalar-valued function of one variable.
  \item $g:\R^2\to\R$ is a scalar-valued function of two variables.
  \item $T$ is a scalar-valued function of three variables. Its natural domain is $z>0$ because $\ln z$ is defined only for positive $z$.
  \item $r:\R\to\R^3$ is a vector-valued function of one variable. It parameterizes a curve in space.
  \item $F:\R^2\to\R^2$ is a vector field on the plane.
\end{enumerate}
\end{examplebox}

\begin{examplebox}[title={2. Reading a level curve}]
Let
\[
  f(x,y)=x^2+2y^2.
\]
Describe the level curve $f(x,y)=8$.

\textbf{Solution.}
The level curve is
\[
  x^2+2y^2=8.
\]
Dividing by $8$ gives
\[
  \frac{x^2}{8}+\frac{y^2}{4}=1.
\]
This is an ellipse centered at the origin. The $x$-intercepts are $\pm\sqrt8=\pm2\sqrt2$, and the $y$-intercepts are $\pm2$.

\textbf{Interpretation.} On this ellipse, every point has the same function value $8$. If $f$ is a loss function, all parameter pairs on this ellipse have equal loss. If $f$ is height, all points on this ellipse have equal elevation.
\end{examplebox}

\begin{examplebox}[title={3. From a table to partial rates of change}]
Suppose $T(x,y)$ is temperature on a plate. The table below gives approximate values.

\begin{center}
\begin{tabular}{c|ccc}
\toprule
$T(x,y)$ & $y=0$ & $y=1$ & $y=2$\\
\midrule
$x=0$ & 80 & 78 & 72\\
$x=1$ & 76 & 74 & 68\\
$x=2$ & 64 & 62 & 56\\
\bottomrule
\end{tabular}
\end{center}

Estimate the rate of change in the $x$ direction near $(1,1)$ using centered differences.

\textbf{Solution.}
Holding $y=1$ fixed, use the values at $x=0$ and $x=2$:
\[
  \frac{T(2,1)-T(0,1)}{2-0}
  =\frac{62-78}{2}
  =-8.
\]
So near $(1,1)$, the temperature decreases by about $8$ units per unit increase in $x$, with $y$ held fixed.

\textbf{Interpretation.} This is the numerical idea behind a partial derivative. Instead of needing a formula, we estimate local change from nearby data.
\end{examplebox}

\begin{examplebox}[title={4. A loss function as a multivariable function}]
Let one data point be $(x,y)=(2,5)$, and consider the prediction model
\[
  \widehat y=wx+b.
\]
The squared error for this one point is
\[
  L(w,b)=(2w+b-5)^2.
\]
Find all parameter pairs $(w,b)$ with zero loss.

\textbf{Solution.}
Zero loss means
\[
  (2w+b-5)^2=0.
\]
Thus
\[
  2w+b-5=0,
\]
or
\[
  b=5-2w.
\]
So all points $(w,b)$ on the line $b=5-2w$ give perfect prediction for this single data point.

\textbf{Interpretation.} A single data point does not determine a unique line. With more data points, the loss surface changes, and the best parameters usually become more constrained.
\end{examplebox}

\section{Python mini-lab: from curves to surfaces to loss landscapes}

This mini-lab previews the computational style of the book. Each full chapter can also have a separate notebook in the \texttt{labs/} folder.

\subsection*{Task A. A one-variable curve}
\begin{lstlisting}[style=pythonstyle]
import numpy as np
import matplotlib.pyplot as plt

x = np.linspace(-3, 3, 400)
y = x**3 - 3*x

plt.figure()
plt.plot(x, y)
plt.axhline(0, linewidth=0.8)
plt.axvline(0, linewidth=0.8)
plt.xlabel("x")
plt.ylabel("f(x)")
plt.title("One-variable function: f(x)=x^3-3x")
plt.show()
\end{lstlisting}

\subsection*{Task B. A two-variable contour plot}
\begin{lstlisting}[style=pythonstyle]
x = np.linspace(-3, 3, 250)
y = np.linspace(-3, 3, 250)
X, Y = np.meshgrid(x, y)
Z = X**2 - Y**2

plt.figure()
cs = plt.contour(X, Y, Z, levels=21)
plt.clabel(cs, inline=True, fontsize=8)
plt.xlabel("x")
plt.ylabel("y")
plt.gca().set_aspect("equal")
plt.title("Level curves of f(x,y)=x^2-y^2")
plt.show()
\end{lstlisting}

\subsection*{Task C. A first numerical gradient}
For
\[
  f(x,y)=x^2+2y^2,
\]
we know exactly that
\[
  \grad f(x,y)=\vect{2x,4y}.
\]
The following code approximates the gradient at a point.
\begin{lstlisting}[style=pythonstyle]
def f_xy(x, y):
    return x**2 + 2*y**2

def numerical_gradient(f, x0, y0, h=1e-5):
    df_dx = (f(x0 + h, y0) - f(x0 - h, y0)) / (2*h)
    df_dy = (f(x0, y0 + h) - f(x0, y0 - h)) / (2*h)
    return np.array([df_dx, df_dy])

point = (1.0, 1.0)
print("Numerical gradient:", numerical_gradient(f_xy, *point))
print("Exact gradient:    ", np.array([2*point[0], 4*point[1]]))
\end{lstlisting}

\subsection*{Task D. A first gradient descent experiment}
\begin{lstlisting}[style=pythonstyle]
def f_val(x, y):
    return x**2 + 2*y**2

def grad_f(x, y):
    return np.array([2*x, 4*y])

eta = 0.12
point = np.array([2.5, 2.0])
path = [point.copy()]

for _ in range(25):
    point = point - eta * grad_f(point[0], point[1])
    path.append(point.copy())

path = np.array(path)

x = np.linspace(-3, 3, 250)
y = np.linspace(-3, 3, 250)
X, Y = np.meshgrid(x, y)
Z = f_val(X, Y)

plt.figure()
plt.contour(X, Y, Z, levels=20)
plt.plot(path[:, 0], path[:, 1], marker="o")
plt.xlabel("x")
plt.ylabel("y")
plt.gca().set_aspect("equal")
plt.title("Gradient descent for f(x,y)=x^2+2y^2")
plt.show()
\end{lstlisting}

\section{Skills checklist}

By the end of this chapter, you should be able to do the following.

\begin{center}
\begin{tabularx}{0.85\textwidth}{X c}
\toprule
\textbf{Skill} & \textbf{Can you do it?}\\
\midrule
Identify the domain and codomain of a function & $\square$\\
Distinguish scalar-valued and vector-valued functions & $\square$\\
Explain why $f(x,y)$ has a surface graph & $\square$\\
Interpret a level curve $f(x,y)=c$ & $\square$\\
Explain why multivariable derivatives depend on direction & $\square$\\
Connect gradients to steepest increase and optimization & $\square$\\
Connect double integrals to volume, mass, and probability & $\square$\\
Interpret a loss function as a function on parameter space & $\square$\\
Use Python to plot curves, contours, surfaces, and simple gradients & $\square$\\
\bottomrule
\end{tabularx}
\end{center}

\section{Practice problems}

\subsection*{Conceptual problems}
\begin{enumerate}
  \item Give three examples of functions of one variable and three examples of functions of several variables. For each example, identify the input and output.
  \item Explain the difference between a point in physical space and a point in parameter space.
  \item Why can $f:\R^3\to\R$ not be graphed in ordinary three-dimensional space in the same way as $f:\R^2\to\R$?
  \item Explain why a contour map is useful even when a surface plot is available.
  \item In your own words, explain why derivatives become more complicated when the input has more than one dimension.
\end{enumerate}

\subsection*{Computation problems}
\begin{enumerate}[resume]
  \item Classify each function by domain and codomain:
  \begin{enumerate}
    \item $f(x)=e^{-x^2}$
    \item $g(x,y)=x^2-y^2$
    \item $h(x,y,z)=\sqrt{x^2+y^2+z^2}$
    \item $r(t)=\vect{t,t^2,t^3}$
    \item $F(x,y)=\vect{x+y,x-y}$
  \end{enumerate}
  \item For $f(x,y)=x^2+y^2$, describe the level curves $f(x,y)=1$, $f(x,y)=4$, and $f(x,y)=9$.
  \item For $f(x,y)=x^2+4y^2$, describe the level curve $f(x,y)=16$.
  \item Compute $\grad f(x,y)$ for $f(x,y)=3x^2+xy+y^2$.
  \item Compute $\grad T(x,y,z)$ for $T(x,y,z)=x^2e^y+\ln z$.
  \item Let $L(w,b)=(2w+b-5)^2+(w+b-3)^2$. Compute $L(1,1)$ and $L(2,1)$.
  \item Let $p(x,y)=Ce^{-(x^2+y^2)}$. Explain what condition $C$ must satisfy if $p$ is to be a probability density on $\R^2$.
\end{enumerate}

\subsection*{Python problems}
\begin{enumerate}[resume]
  \item Modify the contour plot for $f(x,y)=x^2+y^2$ to plot $f(x,y)=x^2+4y^2$. Describe how the level curves change.
  \item Write a Python function for $f(x,y)=\sin(xy)$. Plot its contour map on $[-3,3]\times[-3,3]$.
  \item Numerically approximate the gradient of $f(x,y)=\sin x+\cos y$ at $(0,0)$ using finite differences.
  \item Create a least-squares loss surface for data generated from $y=2x+1$ with small random noise.
  \item Run gradient descent on $f(x,y)=x^2+y^2$ from three different starting points. Plot all three paths.
\end{enumerate}

\subsection*{Discussion and writing problems}
\begin{enumerate}[resume]
  \item Write a short paragraph explaining the phrase: ``A loss function is a landscape.''
  \item Write a short paragraph explaining why multiple integrals are important for probability and statistics.
  \item Choose one application area: physics, economics, biology, statistics, AI, or machine learning. Describe one natural function of several variables in that area.
\end{enumerate}

\section{Chapter summary}

Single-variable calculus studies functions with one input. Multivariable calculus studies functions whose inputs live in spaces such as $\R^2$, $\R^3$, and $\R^n$. This makes the geometry richer and the applications broader.

The main ideas of the book are already visible:
\begin{itemize}
  \item A function $f:\R^n\to\R$ can represent a surface, density, loss, risk, energy, or potential.
  \item When there are many input directions, derivatives become partial derivatives, directional derivatives, gradients, Jacobians, and Hessians.
  \item Integrals become accumulation over regions, solids, surfaces, and probability spaces.
  \item Optimization becomes the study of landscapes in parameter space.
  \item Python helps us visualize, approximate, and experiment with objects that are difficult or impossible to draw by hand.
\end{itemize}

The next chapters build the geometric foundation: Euclidean space, vectors, dot products, projections, lines, planes, cross products, and curves. These objects are the language needed for functions, derivatives, integrals, and vector calculus.

 

\end{document}
